% GENERATED FILE. Edit canonical lessons, then run npm run generate:books.
% canonical-source-hash: fnv1a64:869a2a5ed96f0517
% canonical-lessons: TE-C50-now, TE-C50-day-after, TE-C50-journey, TE-C50-set-out, TE-C50-farewell

\chapter{Taking Your Leave}
\label{ch:te-leave-more}
\hlchaptermodality{50}

\begin{chapteropening}
\textbf{By the end of this chapter:} \emph{I can take my leave in Telugu the way a Telugu speaker actually does, and say why the flat \textquotedblleft{}I am going\textquotedblright{} is the one thing never said at a door.}

Complete TE-C50-farewell and demonstrate the chapter promise.
\end{chapteropening}

\section[ippuḍu]{\te{ఇప్పుడు} (ippuḍu) — now}

\label{lesson:TE-C50-now}



\begin{quote}
Before the new run: what did \emph{alāgē} mean?
\end{quote}

\subsection*{You'll want to know: \te{ఇప్పుడు}}
\textbf{\te{ఇప్పుడు}} (\emph{ippuḍu}) — \textquotedblleft{}now\textquotedblright{}.

The grid from the last chapter has a time row, and this is its near corner. \textbf{-\te{ప్పుడు}} is the time-piece; put a pointing vowel in front of it and the word makes itself.

\te{ఇప్పుడు} is \textquotedblleft{}this time\textquotedblright{} — now. \te{అప్పుడు} is \textquotedblleft{}that time\textquotedblright{} — then. \te{ఎప్పుడు} is \textquotedblleft{}which time\textquotedblright{} — when? You already owned the vowels; the row costs you one word.

\begin{grammarlens}[title={What you've built}]
Now — and a shape that gives you two more words for free.
\end{grammarlens}

\subsection*{Your turn}
Say these aloud:

\begin{itemize}
\raggedright
  \item \emph{ippuḍu}
  \item it once more, slowly
  \item \emph{ippuḍu}, then \emph{ippuḍu veḷḷu} — \textquotedblleft{}go now\textquotedblright{}
\end{itemize}

\begin{tcolorbox}[breakable,colback=teal!4,colframe=teal!35!black,title={Before you move on}]
What does \te{ఇప్పుడు} mean? (\textquotedblleft{}now\textquotedblright{}.) Say it once more.
\end{tcolorbox}
\section[elluṇḍi]{\te{ఎల్లుండి} (elluṇḍi) — the day after tomorrow}

\label{lesson:TE-C50-day-after}



\begin{quote}
Before the new one: what did \emph{ippuḍu} mean?
\end{quote}

\subsection*{You'll want to know: \te{ఎల్లుండి}}
\textbf{\te{ఎల్లుండి}} (\emph{elluṇḍi}) — \textquotedblleft{}the day after tomorrow\textquotedblright{}.

One settled noun for the day after tomorrow. English needs four words and still sounds unsure of itself; Telugu, like its sisters, simply has the word.

There is an older tomorrow hiding at the front of it. \te{ఎల్లి} once meant tomorrow, and though \te{రేపు} took that job in everyday speech, \te{ఎల్లి} is still in service inside \te{ఎల్లుండి}, doing its old work one day further out.

\begin{grammarlens}[title={What you've built}]
Two: now, and the day after tomorrow.
\end{grammarlens}

\subsection*{Your turn}
Say these aloud:

\begin{itemize}
\raggedright
  \item \emph{elluṇḍi}
  \item it once more, slowly
  \item \emph{ippuḍu}, then \emph{elluṇḍi}, and say which one is further off
\end{itemize}

\begin{tcolorbox}[breakable,colback=teal!4,colframe=teal!35!black,title={Before you move on}]
What does \te{ఎల్లుండి} mean? (\textquotedblleft{}the day after tomorrow\textquotedblright{}.) Say it once more.
\end{tcolorbox}
\section[prayāṇaṁ]{\te{ప్రయాణం} (prayāṇaṁ) — a journey}

\label{lesson:TE-C50-journey}



\begin{quote}
Before the new one: what did \emph{elluṇḍi} mean?
\end{quote}

\subsection*{You'll want to know: \te{ప్రయాణం}}
\textbf{\te{ప్రయాణం}} (\emph{prayāṇaṁ}) — \textquotedblleft{}a journey\textquotedblright{}.

Sanskrit \emph{pra-yāṇa}: the root \emph{yā}, \textquotedblleft{}to go\textquotedblright{}, with \emph{pra-}, \textquotedblleft{}forth\textquotedblright{}. A journey is a going-forth, and the word is transparent to anyone holding the pieces.

The ending is the \textbf{-\te{ం}} you have now sorted three times: the tatsama neuter, as on \te{దీపం}, \te{హృదయం} and \te{భోజనం}. When a Telugu noun ends in \textbf{-\te{ం}} and names an abstraction, Sanskrit is usually behind it.

\begin{grammarlens}[title={What you've built}]
Three.
\end{grammarlens}

\subsection*{Your turn}
Say these aloud:

\begin{itemize}
\raggedright
  \item \emph{prayāṇaṁ}
  \item it once more, slowly
  \item \emph{prayāṇaṁ}, then \emph{dīpaṁ}, and name the ending they share
\end{itemize}

\begin{tcolorbox}[breakable,colback=teal!4,colframe=teal!35!black,title={Before you move on}]
What does \te{ప్రయాణం} mean? (\textquotedblleft{}a journey\textquotedblright{}.) Say it once more.
\end{tcolorbox}
\section[bayaludēru]{\te{బయలుదేరు} (bayaludēru) — to set out, to start off}

\label{lesson:TE-C50-set-out}



\begin{quote}
Before the new one: what did \emph{prayāṇaṁ} mean?
\end{quote}

\subsection*{You'll want to know: \te{బయలుదేరు}}
\textbf{\te{బయలుదేరు}} (\emph{bayaludēru}) — \textquotedblleft{}to set out, to start off\textquotedblright{}.

Native, and a picture rather than an abstraction. \te{బయలు} is open ground, the outside; \te{దేరు} is to rise, to come forth. To set out is to come out into the open.

Set it beside \te{ప్రయాణం} and you have two languages building the same idea from opposite ends: Sanskrit from the going, Telugu from the doorstep.

\begin{grammarlens}[title={What you've built}]
Four.
\end{grammarlens}

\subsection*{Your turn}
Say these aloud:

\begin{itemize}
\raggedright
  \item \emph{bayaludēru}
  \item it once more, slowly
  \item \emph{bayaludēru}, then \emph{prayāṇaṁ}, and say which of the two is inherited
\end{itemize}

\begin{tcolorbox}[breakable,colback=teal!4,colframe=teal!35!black,title={Before you move on}]
What does \te{బయలుదేరు} mean? (\textquotedblleft{}to set out, to start off\textquotedblright{}.) Say it once more.
\end{tcolorbox}
\section[veḷḷostānu]{\te{వెళ్ళొస్తాను} (veḷḷostānu) — I'll go and come back --- the everyday goodbye}

\label{lesson:TE-C50-farewell}



\begin{quote}
Before the new one: what did \emph{bayaludēru} mean?
\end{quote}

\subsection*{You'll want to know: \te{వెళ్ళొస్తాను}}
\textbf{\te{వెళ్ళొస్తాను}} (\emph{veḷḷostānu}) — \textquotedblleft{}I'll go and come back --- the everyday goodbye\textquotedblright{}.

Two verbs you already have, welded at the seam: \te{వెళ్ళి}, \textquotedblleft{}having gone\textquotedblright{}, and \te{వస్తాను}, \textquotedblleft{}I will come\textquotedblright{}. \te{వెళ్ళొస్తాను} is \emph{I'll go and come back}, and it is what is said at the door.

What is not said is the flat \textquotedblleft{}I am going\textquotedblright{}. Leaving without the returning half is what is said of the dead, and the rule holds right across the South. The politeness here is not decoration laid over the grammar; it is the grammar. The reply comes back in the same shape: \te{వెళ్ళి} \te{రండి} — \textquotedblleft{}go, and come\textquotedblright{}.

\begin{grammarlens}[title={What you've built}]
Five: \te{ఇప్పుడు}, \te{ఎల్లుండి}, \te{ప్రయాణం}, \te{బయలుదేరు}, and a goodbye that refuses to be one.
\end{grammarlens}

\subsection*{Your turn}
Say these aloud:

\begin{itemize}
\raggedright
  \item \emph{veḷḷostānu}
  \item it once more, slowly
  \item \emph{veḷḷostānu}, then the reply \emph{veḷḷi raṇḍi}, and say why neither half stands alone
\end{itemize}

\begin{tcolorbox}[breakable,colback=teal!4,colframe=teal!35!black,title={Before you move on}]
What does \te{వెళ్ళొస్తాను} mean? (\textquotedblleft{}I'll go and come back --- the everyday goodbye\textquotedblright{}.) Say it once more.
\end{tcolorbox}
